% !TEX root = Appendix_1.tex
% =========================================================================
% APPENDIX 1 -- MACRO DYNAMICS: COINTEGRATION AND THE VEC FOUNDATION
% Econometrica / Econometric Society supplement format.
% Detailed presentation of Section~2.2 results. Replication: Appendix_1/
% =========================================================================
\documentclass[ecta,nameyear,final,supplement]{econsocart}

\usepackage{bm}
\usepackage{booktabs}
\usepackage{array}
\usepackage{float}
\RequirePackage[colorlinks,citecolor=blue,linkcolor=blue,urlcolor=blue,pagebackref]{hyperref}
\urlstyle{same}

\startlocaldefs
\numberwithin{equation}{section}
\numberwithin{table}{section}

\newtheorem{theorem}{Theorem}[section]
\newtheorem{proposition}[theorem]{Proposition}
\newtheorem{lemma}[theorem]{Lemma}
\newtheorem{corollary}[theorem]{Corollary}
\theoremstyle{definition}
\newtheorem{definition}[theorem]{Definition}
\newtheorem{remark}[theorem]{Remark}

\newcommand{\E}{\mathbb{E}}
\newcommand{\KE}{K^{E}}
\newcommand{\KNE}{K^{NE}}

\renewcommand{\thesection}{S\arabic{section}}
% Compact AMS-style labels: S{sec}.{sub}.{subsub}.{para}; 0 if level unused
\makeatletter
\newcommand{\@appzero}[1]{\ifnum\value{#1}=0 0\else\arabic{#1}\fi}
\newcommand{\@appsec}{S\arabic{section}}
\newcommand{\@appsub}{\@appzero{subsection}}
\newcommand{\@appsubsub}{\@appzero{subsubsection}}
\renewcommand{\thesubsection}{\@appsec.\@appsub.0.0}
\renewcommand{\thesubsubsection}{\@appsec.\@appsub.\@appsubsub.0}
\renewcommand{\theparagraph}{\@appsec.\@appsub.\@appsubsub.\arabic{paragraph}}
\setcounter{secnumdepth}{4}
\@addtoreset{subsection}{section}
\@addtoreset{subsubsection}{subsection}
\@addtoreset{paragraph}{section}
\makeatother

% APA 7-style table notes (Publication Manual, \S7.21)
\newcommand{\tabnote}{\par\addvspace{0.4em}\noindent\emph{Note.}}
\newcommand{\tabsource}{\ \textit{Source:} Authors' own calculations based on
the monthly sample in Section~\ref{sec:data} ($T=858$). Software: MATLAB
R2025b (The MathWorks, 2025); Econometrics Toolbox.}
\endlocaldefs

\begin{document}

\begin{frontmatter}

\title{Supplement to ``Identifying Rational Types in Unknown Environments''\\
Appendix 1: Macro Dynamics --- Cointegration and the VEC Foundation}
\runtitle{Macro Dynamics Supplement}

\begin{aug}
\author[add1,add2]{\fnms{Humberto}~\snm{Bernal}}
\address[add1]{%
\orgdiv{Economics Program},
\orgname{Universidad Colegio Mayor de Cundinamarca}}
\address[add2]{%
\orgdiv{Ph.D. in Economics},
\orgname{Universidad de los Andes}}
\ead[label=e1]{hbernal@uniandes.edu.co}
\ead[label=e2]{hbernalc@universidadmayor.edu.co}
\end{aug}

\begin{funding}
Economist, professor of the Economics Program at the Universidad Colegio
Mayor de Cundinamarca and graduate of the Ph.D.\ in Economics from the
Universidad de los Andes, Colombia.
\end{funding}

\end{frontmatter}

\noindent
This supplement presents in full detail the macroeconometric results
summarized in Section~2.2 of the main paper (``Macro dynamics: cointegration
and the VEC foundation''). It documents the data, the
integration-order tests, the Johansen cointegration analysis, the estimated
Vector Error Correction (VEC) model with its adjustment dynamics and residual
diagnostics, the Gonzalo--Granger permanent--transitory decomposition, the
univariate Local Linear Trend (LLT) state-space model, and the empirical
verification of the stochastic-trend equivalence corollary. All estimates
were reproduced from the MATLAB scripts in the \texttt{Appendix\_1/}
replication package (Section~\ref{sec:repro}).

% ============================================================================
\section{Data and sample}
\label{sec:data}
% ============================================================================

The vector of interest is the monthly triple
\begin{equation}
Y_t = \bigl(\KE_t,\; \KNE_t,\; r_t\bigr)',
\label{eq:vector}
\end{equation}
where $\KE_t$ is the number of kidnappings for ransom (extortive kidnapping),
$\KNE_t$ is the number of non-extortive kidnappings, and $r_t$ is the
commercial lending interest rate. Kidnapping counts come from the National
Center of Historical Memory \citep{CNMH2013,cnmh_datos} and official
statistics of the Colombian State \citep{mindefensa_stats}; the interest rate
is the active commercial rate used as the indicator of macro-financial
conditions \citep{Masih1995,Suarez_2025}.

The estimation sample is the continuous monthly series from July~1954 to
December~2025 ($T=858$ observations). All macroeconometric models in this
supplement---unit-root tests, Johansen analysis, the VEC, the LLT, the
Gonzalo--Granger trends, and the trend-equivalence validation of
$d_t$---use this same sample, as implemented in \path{VEC\_2T.m},
\path{LLT\_Extorsivo.m}, and \path{check_trend_equivalence.m}. Macro trend
figures in the main paper (Section~2.2) plot the same July~1954--December~2025
window. Descriptive institutional figures elsewhere in the main text display
1964--2025. Isolated missing values were imputed with a centered moving median
of order three before estimation.

Economically, $r_t$ proxies the macroeconomic environment that fixes the
opportunity cost and the continuation value of crime, while $\KNE_t$ marks
the phase in which the criminal organization accumulates logistic capacity,
operational experience, and territorial control, lowering its marginal cost
of operation and enabling the more systematic organization of extortive
kidnapping.

% ============================================================================
\section{Integration order}
\label{sec:unitroot}
% ============================================================================

Table~\ref{tab:unitroot} reports the Augmented Dickey--Fuller test of
\citet{DickeyFuller1979} and the KPSS test of \citet{Kwiatkowski1992}, in
levels and in first differences, with twelve augmentation lags and a
constant. At the 5\% level, the ADF test does not reject the unit-root null in
levels for any series ($p>0.05$), while the KPSS test rejects stationarity
in levels at the 1\% level ($p\le 0.01$). In first differences, the ADF test
rejects the unit-root null at the 1\% level ($p=0.001$) and the KPSS test
does not reject stationarity at the 10\% level ($p\ge 0.10$). These four
criteria jointly classify $\KE_t$, $\KNE_t$, and $r_t$ as $I(1)$ at the 5\%
level.

\begin{table}[t]
\centering
\caption{Augmented Dickey--Fuller and KPSS Unit-Root Test $p$-Values}
\label{tab:unitroot}
\begin{tabular}{lcccc}
\toprule
& \multicolumn{2}{c}{Levels} & \multicolumn{2}{c}{First differences} \\
\cmidrule(lr){2-3}\cmidrule(lr){4-5}
Series & ADF $p$ & KPSS $p$ & ADF $p$ & KPSS $p$ \\
\midrule
$\KE_t$  & $0.28$ & $<0.01$ & $<0.01$ & $>0.10$ \\
$\KNE_t$ & $0.41$ & $<0.01$ & $<0.01$ & $>0.10$ \\
$r_t$    & $0.44$ & $<0.01$ & $<0.01$ & $>0.10$ \\
\bottomrule
\end{tabular}
\tabnote\ ADF: constant and 12 augmentation lags (\texttt{model = ARD}).
KPSS: level-stationarity test, 12 lags, no trend. MATLAB truncates
$p$-values at $[0.01,\,0.10]$ (KPSS) and $[0.001,\,0.999]$ (ADF). In levels,
ADF does not reject the unit-root null at 5\%; KPSS rejects stationarity at
1\%. In first differences, ADF rejects the unit-root null at 1\%; KPSS does
not reject stationarity at 10\%.
\tabsource
\end{table}

% ============================================================================
\section{Johansen cointegration analysis}
\label{sec:johansen}
% ============================================================================

The cointegration rank is determined with the Johansen trace test
\citep{Johansen1988} on the levels VAR with $p=36$ lags and an unrestricted
constant (MATLAB deterministic specification \texttt{H1}). The lag order was
selected so that the implied VEC$(35)$ exhibits no residual autocorrelation
(Section~\ref{sec:diagnostics}) while remaining parsimonious under the AIC
among admissible specifications. Table~\ref{tab:johansen} reports the
results: the null $r=0$ is rejected at the 1\% level ($p<0.01$), whereas
$r\le 1$ is not rejected at the 5\% level ($p=0.11$) and $r\le 2$ is not
rejected at the 5\% level ($p=0.08$) although it would be rejected at the
10\% level. The evidence therefore supports exactly one cointegrating vector,
$r=1$, with $k-r=2$ common stochastic trends.

\begin{table}[H]
\centering
\caption{Johansen Trace Test for Cointegration Rank}
\label{tab:johansen}
\begin{tabular}{lcccl}
\toprule
Null hypothesis & Trace statistic & 95\% critical value & $p$-value &
Decision \\
\midrule
$r=0$     & $40.32$ & $29.80$ & $<0.01$ & Reject at 1\% \\
$r\le 1$  & $13.17$ & $15.49$ & $0.11$  & Not reject at 5\% \\
$r\le 2$  & $3.02$  & $3.84$  & $0.08$  & Not reject at 5\% \\
\bottomrule
\end{tabular}
\tabnote\ VAR in levels with $p=36$ lags and unrestricted constant
(MATLAB specification \texttt{H1}). Decisions follow the asymptotic
$p$-value at the stated significance level. The null $r\le 2$ would be
rejected at 10\% ($p=0.08$); rank $r=1$ is selected using the 5\%
criterion.
\tabsource
\end{table}

% ============================================================================
\section{The VEC model}
\label{sec:vec}
% ============================================================================

\subsection{Long-run equilibrium}

With $r=1$ and $q=p-1=35$ short-run lag matrices, the estimated VEC is
\begin{equation}
\Delta Y_t = \alpha\,\beta' Y_{t-1}
 + \sum_{j=1}^{35}\Gamma_j\,\Delta Y_{t-j} + c + \varepsilon_t ,
\label{eq:vec-general}
\end{equation}
estimated by maximum likelihood (Johansen reduced-rank regression). The fit
statistics are $\log L = -7{,}667.23$, $\mathrm{AIC} = 15{,}980.45$, and
$\mathrm{BIC} = 17{,}516.19$. The unnormalized cointegrating vector is
\begin{equation}
0.04\,\KE_t \;-\; 0.07\,\KNE_t \;-\; 0.05\,r_t
 \;=\; \tilde\varepsilon_t ,
\qquad \tilde\varepsilon_t \sim I(0),
\label{eq:coint-raw}
\end{equation}
and an ADF test on the error-correction term
$\tilde\varepsilon_t = \beta' Y_t$ rejects the unit-root null at the 1\%
level ($p<0.01$), confirming the stationarity of the long-run relation.
Normalizing \eqref{eq:coint-raw} on $\KE_t$ yields the structural long-run
equation
\begin{equation}
\KE_t \;=\; 1.93\,\KNE_t \;+\; 1.31\,r_t \;+\; \varepsilon_t ,
\qquad \varepsilon_t \sim I(0).
\label{eq:coint-normalized}
\end{equation}

The coefficient on $\KNE_t$ indicates that non-extortive kidnapping operates
as a signal of accumulated coercive and organizational capacity: in the long
run, a unit increase in this modality is associated with an increase of
approximately $1.93$ units in extortive kidnapping, as the expansion of
logistic infrastructure, territorial control, and intimidation lowers the
marginal cost of moving toward higher-rent activities. The coefficient on
the interest rate reflects a liquidity-constraint mechanism: periods of
monetary contraction raise the financing cost of criminal organizations and
push them toward extortive kidnapping as a source of immediate rent.

\subsection{Adjustment dynamics}

Table~\ref{tab:alpha} reports the estimated adjustment (loading) vector
$\alpha$ with standard errors. Written equation by equation, the
error-correction system is
\begin{align}
\Delta \KE_t  &= -2.94\,\tilde\varepsilon_{t-1}
 + \textstyle\sum_{j=1}^{35}\Gamma_{1,j}\,\Delta Y_{t-j} + \varepsilon_{1,t},
\label{eq:vec1}\\
\Delta \KNE_t &= \phantom{-}0.36\,\tilde\varepsilon_{t-1}
 + \textstyle\sum_{j=1}^{35}\Gamma_{2,j}\,\Delta Y_{t-j} + \varepsilon_{2,t},
\label{eq:vec2}\\
\Delta r_t    &= -0.06\,\tilde\varepsilon_{t-1}
 + \textstyle\sum_{j=1}^{35}\Gamma_{3,j}\,\Delta Y_{t-j} + \varepsilon_{3,t}.
\label{eq:vec3}
\end{align}
Only the loading in the $\KE$ equation is statistically significant at the
1\% level ($p<0.01$). The loading in the $\KNE$ equation is not significant
at the 5\% or 10\% level ($p=0.17$). The loading in the $r$ equation is not
significant at the 5\% level ($p=0.10$) but is marginally significant at the
10\% level (exact $p=0.098$). Extortive kidnapping is therefore the variable
that absorbs disequilibria at conventional levels, while $\KNE_t$ and $r_t$
are weakly exogenous with respect to the long-run relation at the 5\% level.
Because $\beta_{\KE} = 0.04$, the effective
adjustment speed is
\begin{equation}
\bigl|\alpha_{\KE}\cdot\beta_{\KE}\bigr| \;=\;
\lvert -2.94 \times 0.04\rvert \;\approx\; 0.11 :
\label{eq:speed}
\end{equation}
approximately $11\%$ of any deviation from the long-run equilibrium is
corrected each month through adjustments in extortive kidnapping.

\begin{table}[t]
\centering
\caption{Vector Error Correction Loading Coefficients}
\label{tab:alpha}
\begin{tabular}{lcccl}
\toprule
Equation & $\hat\alpha$ & Std.\ error & $p$-value & Significance \\
\midrule
$\Delta \KE_t$  & $-2.94$ & $0.67$ & $<0.01$ & 1\% \\
$\Delta \KNE_t$ & $0.36$  & $0.26$ & $0.17$  & --- \\
$\Delta r_t$    & $-0.06$ & $0.04$ & $0.10$  & 10\% \\
\bottomrule
\end{tabular}
\tabnote\ VEC($p=36$, $r=1$), unrestricted constant, maximum-likelihood
estimation ($T=858$). Two-sided $p$-values. The Significance column reports
the highest conventional level at which the coefficient differs from zero
(1\%, 5\%, 10\%); ``---'' denotes not significant at 10\%. The exact
$p$-value for $\Delta r_t$ is $0.098$.
\tabsource
\end{table}

\subsection{Residual diagnostics}
\label{sec:diagnostics}

Table~\ref{tab:diagnostics} summarizes the residual diagnostics of the
estimated system. Ljung--Box tests at twenty lags do not reject the absence
of autocorrelation at any conventional level ($p\ge 0.99$), validating the
dynamic specification with $p=36$. Engle ARCH tests at five lags reject
homoskedasticity at the 1\% level in the $\KE$ and $r$ equations
($p<0.01$) and do not reject it at the 5\% level in the $\KNE$ equation
($p=0.95$). This pattern is expected in long monthly series spanning regimes
of very different volatility (armed-conflict escalation and de-escalation,
interest-rate liberalization) and does not bias the point estimates of
$\beta$ and $\alpha$, although inference on short-run coefficients should be
read with heteroskedasticity-robust caution.

\begin{table}[t]
\centering
\caption{Residual Diagnostic Tests for the Estimated VEC System}
\label{tab:diagnostics}
\begin{tabular}{lccc}
\toprule
Equation & Ljung--Box$(20)$ $p$ & ARCH$(5)$ $p$ & ARCH decision \\
\midrule
$\Delta \KE_t$  & $1.00$ & $<0.01$ & Reject at 1\% \\
$\Delta \KNE_t$ & $0.99$ & $0.95$  & Not reject at 5\% \\
$\Delta r_t$    & $1.00$ & $<0.01$ & Reject at 1\% \\
\bottomrule
\end{tabular}
\tabnote\ VEC($p=36$, $r=1$), 35 short-run lags. Ljung--Box test on VEC
residuals at 20 lags; Engle ARCH LM test at 5 lags. ARCH decision: reject
homoskedasticity at 1\% or do not reject at 5\%, as indicated.
\tabsource
\end{table}

% ============================================================================
\section{Gonzalo--Granger permanent--transitory decomposition}
\label{sec:gg}
% ============================================================================

With $k=3$ variables and $r=1$ cointegrating vector, the decomposition of
\citet{GonzaloGranger1995} identifies $k-r=2$ common stochastic trends,
$f_t = \alpha_\perp' Y_t$, where $\alpha_\perp$ spans the orthogonal
complement of the adjustment vector. With the rational normalization that
assigns unit weight to $\KNE_t$ and $r_t$ respectively, the estimated
factors are
\begin{align}
f_{1,t} &= 0.12\,\KE_t + 1.00\,\KNE_t ,
\label{eq:gg-f1}\\
f_{2,t} &= -0.02\,\KE_t + 1.00\,r_t .
\label{eq:gg-f2}
\end{align}
The long-run representation of the system,
$Y_t = A\,f_t + \mu_t$ with $A = \beta_\perp(\alpha_\perp'\beta_\perp)^{-1}$
and $\mu_t \sim I(0)$, has estimated loading matrix
\begin{equation}
\begin{pmatrix} \KE_t \\[2pt] \KNE_t \\[2pt] r_t \end{pmatrix}
=
\begin{pmatrix}
1.60 & 1.09 \\
0.81 & -0.13 \\
0.04 & 1.02
\end{pmatrix}
\begin{pmatrix} f_{1,t} \\[2pt] f_{2,t} \end{pmatrix}
+
\begin{pmatrix} \mu_{1,t} \\[2pt] \mu_{2,t} \\[2pt] \mu_{3,t} \end{pmatrix},
\qquad \mu_t \sim I(0).
\label{eq:gg-loading}
\end{equation}

The factor $f_{1,t}$ captures persistent coercive capacity (dominated by
non-extortive kidnapping), while $f_{2,t}$ reflects permanent macro-financial
conditions (dominated by the interest rate). The loadings of $\KE_t$
($1.60$ on $f_1$ and $1.09$ on $f_2$) exceed unity, indicating that
extortive kidnapping amplifies both permanent forces. The long-run trend of
extortive kidnapping implied by the decomposition,
$\tau^{\mathrm{VEC}}_{\KE,t} = 1.60\,f_{1,t} + 1.09\,f_{2,t}$, is the
VEC--GG trend used in the equivalence exercise of
Section~\ref{sec:equivalence}.

% ============================================================================
\section{The Local Linear Trend model}
\label{sec:llt}
% ============================================================================

As an independent univariate contrast, extortive kidnapping is decomposed in
state-space form \citep{Harvey1989} into a stochastic level, a persistent
slope, and a stationary cycle:
\begin{equation}
\begin{aligned}
\text{Measurement:} \quad & y_t = \tau_t + c_t + \xi_t,
&& \xi_t \sim \mathcal{N}(0,\sigma_\xi^2), \\[4pt]
\text{Trend:} \quad & \tau_t = \tau_{t-1} + b_{t-1} + \eta_t,
&& \eta_t \sim \mathcal{N}(0,\sigma_\eta^2), \\[4pt]
\text{Slope:} \quad & b_t = \textstyle\sum_{j=1}^{8}\phi_j\, b_{t-j} + \zeta_t,
&& \zeta_t \sim \mathcal{N}(0,\sigma_\zeta^2), \\[4pt]
\text{Cycle:} \quad & c_t = \textstyle\sum_{k=1}^{5}\rho_k\, c_{t-k} + \kappa_t,
&& \kappa_t \sim \mathcal{N}(0,\sigma_\kappa^2),
\end{aligned}
\label{eq:llt}
\end{equation}
estimated by maximum likelihood via the Kalman filter on the same monthly
sample ($T=858$, $\log L = -3{,}847.34$, $\mathrm{AIC}=7{,}728.68$,
$\mathrm{BIC}=7{,}809.51$; 17 free parameters). Table~\ref{tab:llt} reports
the full parameter set of the replication run. All four innovation standard
deviations are significant at the 1\% level ($p<0.01$ on the log-scale
parameters):
\begin{equation}
\hat\sigma_\eta \approx 7.13, \qquad
\hat\sigma_\zeta \approx 3.51, \qquad
\hat\sigma_\kappa \approx 6.77, \qquad
\hat\sigma_\xi \approx 11.96 .
\label{eq:llt-sigmas}
\end{equation}

\begin{remark}
Quasi-Newton maximum likelihood on this 15-state model can converge to
slightly different points across MATLAB versions and starting values; the
innovation standard deviations in \eqref{eq:llt-sigmas} are from the
replication package run of \path{LLT_Extorsivo.m}. Small numerical
differences do not affect any qualitative conclusion: in every run the
permanent-shock variance $\sigma_\eta^2$ is large and highly significant.
\end{remark}

The significance of $\hat\sigma_\eta$ at the 1\% level confirms that
permanent shocks dominate the long-run level of extortive kidnapping. Among
the slope AR coefficients, $\phi_2$, $\phi_3$, $\phi_5$, and $\phi_7$ are
significant at the 1\% level; $\phi_8$ is significant at the 10\% level
($p=0.06$) but not at the 5\% level; $\phi_1$, $\phi_4$, and $\phi_6$ are not
significant at the 10\% level. All five cycle AR coefficients ($\rho_1$
through $\rho_5$) are significant at the 1\% level. The order-8
autoregressive structure of the slope captures gradual changes in the growth
rate, consistent with slow adjustments in criminal capacity rather than
abrupt breaks. The order-5 stationary cycle absorbs medium-run fluctuations
without distorting the permanent component.

\begin{table}[t]
\centering
\caption{Local Linear Trend State-Space Maximum-Likelihood Estimates}
\label{tab:llt}
\begin{tabular}{lcccl}
\toprule
Parameter & Estimate & Std.\ error & $p$-value & Sig. \\
\midrule
$\log\sigma_\eta$ (level)      & $1.96$  & $0.06$ & $<0.01$ & 1\% \\
$\log\sigma_\zeta$ (slope)     & $1.26$  & $0.20$ & $<0.01$ & 1\% \\
$\phi_1$ (slope AR)            & $-0.38$ & $0.33$ & $0.24$  & --- \\
$\phi_2$                       & $-0.28$ & $0.09$ & $<0.01$ & 1\% \\
$\phi_3$                       & $-0.51$ & $0.16$ & $<0.01$ & 1\% \\
$\phi_4$                       & $-0.02$ & $0.07$ & $0.72$  & --- \\
$\phi_5$                       & $0.50$  & $0.05$ & $<0.01$ & 1\% \\
$\phi_6$                       & $-0.05$ & $0.15$ & $0.73$  & --- \\
$\phi_7$                       & $0.67$  & $0.11$ & $<0.01$ & 1\% \\
$\phi_8$                       & $0.59$  & $0.31$ & $0.06$  & 10\% \\
$\log\sigma_\kappa$ (cycle)    & $1.91$  & $0.08$ & $<0.01$ & 1\% \\
$\rho_1$ (cycle AR)            & $-0.76$ & $0.04$ & $<0.01$ & 1\% \\
$\rho_2$                       & $-0.18$ & $0.02$ & $<0.01$ & 1\% \\
$\rho_3$                       & $-0.24$ & $0.03$ & $<0.01$ & 1\% \\
$\rho_4$                       & $-0.82$ & $0.04$ & $<0.01$ & 1\% \\
$\rho_5$                       & $-0.63$ & $0.04$ & $<0.01$ & 1\% \\
$\log\sigma_\xi$ (irregular)   & $2.48$  & $0.05$ & $<0.01$ & 1\% \\
\bottomrule
\end{tabular}
\tabnote\ Kidnapping-for-ransom series; slope AR(8), cycle AR(5), stochastic
level. Variance parameters estimated in logarithms. The Sig.\ column reports
the highest conventional significance level (1\%, 5\%, 10\%); ``---'' denotes
not significant at 10\%.
\tabsource
\end{table}

% ============================================================================
\section{Equivalence of stochastic trends}
\label{sec:equivalence}
% ============================================================================

The VEC--GG trend of Section~\ref{sec:gg} and the LLT trend of
Section~\ref{sec:llt} are built from conceptually independent principles:
multivariate cointegration restrictions versus univariate state-space
filtering. Their sustained coherence is the empirical manifestation of the
following corollary, stated and proved as
Corollary~A1.2 in the proofs supplement (Appendix~2).

\begin{corollary}[Equivalence of stochastic trends]
\label{cor:equiv-trend}
Let $\{y_t\}$ be a time series integrated of order one, $I(1)$. Suppose the
existence of two valid decompositions defined as
\begin{equation*}
y_t = P_t^{(1)} + u_t^{(1)} = P_t^{(2)} + u_t^{(2)},
\end{equation*}
where $P_t^{(1)}$ and $P_t^{(2)}$ are processes integrated of order one and
$u_t^{(1)}, u_t^{(2)}$ are stationary $I(0)$ processes. Then
\begin{equation*}
P_t^{(1)} - P_t^{(2)} \sim I(0).
\end{equation*}
Consequently, all valid stochastic trends derived from the same observed
series are equivalent up to a stationary component.
\end{corollary}

Empirically, the difference
$d_t = \tau^{\mathrm{VEC}}_{\KE,t} - \tau^{\mathrm{LLT}}_{\KE,t}$ between
the two estimated trends behaves as the corollary predicts. Trends and
unit-root tests on $d_t$ use the full estimation sample July~1954--December~2025
($T=858$), as in \path{check_trend_equivalence.m}. The ADF test rejects a unit
root in $d_t$ at the
1\% level ($p<0.01$). The KPSS test with a Newey--West bandwidth of twelve
lags does not reject level stationarity at the 5\% level ($p=0.07$) but
would reject it at the 10\% level ($p<0.10$); with MATLAB's short default
bandwidth the KPSS test rejects stationarity at the 5\% level ($p=0.01$),
reflecting the persistence of $d_t$ at short horizons rather than a stochastic
trend. Following the corollary, the ADF result at the 1\% level together
with the KPSS(12) result at the 5\% level support $d_t\sim I(0)$: the two
estimates do not diverge systematically in the long run. The VEC through the
Gonzalo--Granger decomposition and the LLT through agnostic state-space
filtering independently reveal the same permanent component of the
phenomenon, which makes the long-run characterization robust to the choice
of econometric framework.

% ============================================================================
\section{Relevance for the mechanism}
\label{sec:relevance}
% ============================================================================

\paragraph{Scope.}
This section explains how the macroeconometric results in
Sections~\ref{sec:data}--\ref{sec:equivalence} feed the mechanism-design model
of the main paper (Section~2.2, ``Macro dynamics: cointegration and the VEC
foundation''). It is \emph{interpretive}: it introduces no new estimates. Every
empirical quantity cited below is reported in Sections~S2--S7 and can be
verified from the downloaded replication package using
\path{README.txt} at the package root---without access to the author's
working directories or any path outside the folder delivered to the reader.

\paragraph{Long-run equilibrium and the interest-rate channel.}
The long-run equilibrium \eqref{eq:coint-normalized} establishes that the
interest rate is a structural determinant of extortive kidnapping: monetary
contraction and credit restriction raise the cost of capital for criminal
organizations, reshaping their incentives toward extraction strategies with
faster returns ($\hat\beta_r/\hat\beta_{\KE}\approx 1.31$ in the normalized
relation). This macroeconomic environment is not neutral with respect to
kidnapper types: it defines the support of the cost distribution and of the
ransom valuation faced by heterogeneous agents in each period, feeding the
parameter space within which the screening mechanism of the main paper
operates.

\paragraph{Common trends and type heterogeneity.}
The persistence of the two common trends identified in Section~\ref{sec:gg}
(coercive capacity $f_{1,t}$ and macro-financial conditions $f_{2,t}$)
implies that heterogeneity across types is neither transitory nor purely
idiosyncratic, but structurally conditioned by the long-run macroeconomic
environment; this property justifies its explicit treatment as a restriction
on the type support in the mechanism-design model.

\paragraph{Trend equivalence and calibration anchors.}
The trend-equivalence result of Section~\ref{sec:equivalence} closes the
model-dependence argument: the permanent component is unique up to stationary
perturbations (Corollary~\ref{cor:equiv-trend}, proved in Appendix~2), so the
long-run parameters carried into the calibration are structural anchors rather
than artifacts of one econometric framework.

\paragraph{Reading guide for replicators.}
Table~\ref{tab:mechanism-map} maps each mechanism-relevant input to its
location in this supplement and in the replication package. A reader who
downloads the package can proceed as follows.
\begin{enumerate}
\item Read \path{tex/Appendix_1.pdf} (or this \LaTeX{} source) through
Section~S7, then Section~S8.
\item Open \path{README.txt} for folder layout and software requirements.
\item To \emph{rerun} all estimates: execute Steps~1--3 in \path{README.txt}
(MATLAB R2025b, Econometrics Toolbox).
\item To \emph{inspect outputs without rerunning MATLAB}: open the pre-exported
CSV files listed in the last column of Table~\ref{tab:mechanism-map}; trend
equivalence can be read from \path{matlab/Trend_Difference/trend_difference.csv}
(column \texttt{Difference}).
\end{enumerate}
Section~\ref{sec:repro} lists every script; the formal mechanism proofs and
microeconometric priors ($\mu_0$, hazards) are in the main paper and its
other supplements, not in this package.

\begin{table}[t]
\centering
\caption{Mapping of Mechanism Inputs to Replicated Macro Results}
\label{tab:mechanism-map}
\begin{tabular}{p{0.22\linewidth}p{0.30\linewidth}p{0.38\linewidth}}
\toprule
Mechanism input & Where in this supplement & Replication package \\
\midrule
$I(1)$ order of $\KE_t$, $\KNE_t$, $r_t$
  & Section~\ref{sec:unitroot}, Table~\ref{tab:unitroot}
  & \path{matlab/VEC/VEC\_2T.m} (Step~1 in \path{README.txt}) \\[4pt]
Cointegration rank $r=1$
  & Section~\ref{sec:johansen}, Table~\ref{tab:johansen}
  & \path{matlab/VEC/VEC\_2T.m} \\[4pt]
Long-run relation
  $\KE_t = 1.93\,\KNE_t + 1.31\,r_t + \varepsilon_t$
  & \eqref{eq:coint-normalized}, Section~\ref{sec:vec}
  & \path{matlab/VEC/VEC\_2T.m}; $\beta$ in console output \\[4pt]
Adjustment: only $\KE_t$ loads on the ECT (5\% level)
  & Section~\ref{sec:vec}, Table~\ref{tab:alpha}
  & \path{matlab/VEC/VEC\_2T.m} \\[4pt]
Common trends $f_{1,t}$, $f_{2,t}$ and VEC--GG loadings
  & Section~\ref{sec:gg}, \eqref{eq:gg-f1}--\eqref{eq:gg-loading}
  & \path{matlab/VEC/VEC\_2T.m};
    \path{matlab/VEC/Tendencias\_SE.csv} \\[4pt]
LLT permanent component of $\KE_t$
  & Section~\ref{sec:llt}, \eqref{eq:llt-sigmas}, Table~\ref{tab:llt}
  & \path{matlab/SE\_Kalman/LLT\_Extorsivo.m};
    \path{matlab/SE\_Kalman/LLT\_decomposition.csv} \\[4pt]
Trend equivalence: $d_t=\tau^{\mathrm{VEC}}_{\KE}-\tau^{\mathrm{LLT}}_{\KE}
  \sim I(0)$
  & Section~\ref{sec:equivalence}, Corollary~\ref{cor:equiv-trend}
  & \path{matlab/Trend\_Difference/check\_trend\_equivalence.m};
    \path{matlab/Trend\_Difference/trend\_difference.csv} \\[4pt]
Full numeric check vs.\ all tables
  & Sections~S2--S7
  & \path{matlab/Trend\_Difference/verify\_appendix1.m} (optional) \\
\bottomrule
\end{tabular}
\tabnote\ This table is a reading guide only; it does not report new
statistics. CSV files allow inspection of trends and differences without
rerunning MATLAB. The mechanism model, calibration, and microeconometric
priors are documented in the main paper.
\tabsource
\end{table}

% ============================================================================
\section{Reproducibility}
\label{sec:repro}
% ============================================================================

\begin{sloppypar}
All results in this supplement are reproduced by MATLAB scripts in the
\texttt{Appendix\_1/} replication package. No script references paths outside
that folder. See \path{README.txt} at the package root for step-by-step
instructions. Requirements: MATLAB R2025b (or later) with the Econometrics
Toolbox.
\end{sloppypar}

\begin{enumerate}
\item \path{matlab/VEC/VEC_2T.m} --- unit roots, Johansen test, VEC
  estimation, Gonzalo--Granger decomposition, diagnostics, IRFs, and
  forecasts; exports \path{matlab/VEC/Tendencias_SE.csv}.
\item \path{matlab/SE_Kalman/LLT_Extorsivo.m} --- LLT state-space model;
  exports \path{matlab/SE_Kalman/LLT_decomposition.csv}.
\item \path{matlab/Trend_Difference/check_trend_equivalence.m} --- trend
  equivalence (Section~\ref{sec:equivalence}); exports
  \path{matlab/Trend_Difference/trend_difference.csv}.
\item (Optional) \path{matlab/Trend_Difference/verify_appendix1.m} --- full
  numeric check against all tables in this supplement.
\item (Optional) Compile \path{tex/Appendix_1.pdf} from \path{tex/} with
  \texttt{pdflatex} and \texttt{bibtex}.
\end{enumerate}

\begin{sloppypar}
Section~\ref{sec:relevance} (Table~\ref{tab:mechanism-map}) maps macro outputs
to mechanism inputs. Input data: \path{matlab/VEC/Data.csv} ($T=858$).
\end{sloppypar}

\bibliographystyle{econsoc}
\bibliography{references}

\end{document}
